\chapter{热力学第二定律：熵}\label{ch:b1:second-law-entropy}

把一只热勺子放入一杯冷水中：勺子变冷，水变暖，
两者最终达到同一温度。第一定律同样允许
其逆过程——水继续变冷而勺子变热——
能量精确到焦耳都守恒，但这种过程从不发生。
倒放的影片一眼即可识别；分子所遵循的定律
没有方向之分，而世界却有。热力学第二定律
指出了区分两个方向的量：
熵，封闭系统只能产生熵，不能消灭熵。本章
陈述该定律，计算本分册各系统的熵，
说明熵如何度量不可逆性，并在
最后给出它的计数意义。

\begin{omfigure}
\includegraphics[width=0.42\linewidth]{images/book3/ai/b3-23-coffee-ice-cube.jpg}

{\small 冰块在热咖啡中熔化：热量只沿一个方向流动，所产生的熵度量了混合偏离可逆的程度。}
\end{omfigure}

\section{热力学第二定律}

\begin{theorem}[热力学第二定律]\label{thm:b1:second-law-entropy:secondlaw}
\index{热力学第二定律}\index{熵}\index{熵!产生}\index{熵!交换}
每个封闭系统都拥有一个广延状态函数，
\emph{熵} $S$（\unit{J/K}），使得任一过程中
\[
\Delta S = S_{\mathrm{exch}} + S_{\mathrm{created}}, \qquad
S_{\mathrm{exch}} = \int\frac{\delta Q}{T_{\mathrm{ext}}}, \qquad
S_{\mathrm{created}} \geq 0 ,
\]
其中 $\delta Q$ 为从外界接收的热量，
它在进入处的温度 $T_{\mathrm{ext}}$ 下传入，并且
当且仅当过程可逆时 $S_{\mathrm{created}} = 0$。
孤立系统（$\delta Q = 0$）的熵只能
增大，其平衡态即为熵最大的状态。
\end{theorem}
\admitted

\begin{remark}[理解该定律]\label{rem:b1:second-law-entropy:reading}
与能量不同，熵并不守恒：它与热量\emph{交换}
（且只与热量交换，功不携带熵），并由每个不可逆过程\emph{产生}，
如摩擦力、热量穿越温差传递、
气体迅速冲入真空、混合。“产生项”即为不可逆性的度量，
而过程只有在该项非负时才可能发生。
“可逆”是一种理想化——准静态的、
无摩擦的、热量在温差趋于零时交换——
真实过程只能无限接近而永远无法达到。
\end{remark}

\begin{omfigure}
\begin{tikzpicture}[font=\small, x=1cm, y=1cm]
  \draw[thick, rounded corners, fill=omDef!8] (0,0) rectangle (4.2,2.6);
  \node at (2.1,1.9) {系统};
  \node[omDef] at (2.1,1.15) {$\Delta S = S_{\mathrm{exch}} + S_{\mathrm{created}}$};
  \node[omThm, font=\footnotesize] at (2.1,0.5) {$S_{\mathrm{created}} \geq 0$（内部）};
  \draw[-Stealth, omProp, very thick] (-3.6,1.1) -- (-0.1,1.1) node[midway, above] {$\delta Q$ 在 $T_{\mathrm{ext}}$};
  \node[omProp, font=\footnotesize] at (-1.85,0.65) {$S_{\mathrm{exch}} = \int\delta Q/T_{\mathrm{ext}}$};
  \draw[-Stealth, omExo, very thick] (-3.6,2.2) -- (-0.1,2.2) node[midway, above] {$W$（无熵）};
\end{tikzpicture}

{\small 熵平衡：热量携带熵 $\delta Q/T_{\mathrm{ext}}$
穿过边界，功不携带熵，而系统内部的不可逆过程
产生部分熵——但绝不消灭它。}
\end{omfigure}

\begin{proposition}[基本恒等式；理想气体与凝聚相的熵]\label{prop:b1:second-law-entropy:identity}
\index{热力学基本恒等式}\index{熵!理想气体}
对组成固定的封闭系统，沿任一可逆路径，
$\delta Q_{\mathrm{rev}} = T\,\dd S$ 且 $\delta W_{\mathrm{rev}} = -P\,\dd V$，于是
\[
\dd U = T\,\dd S - P\,\dd V ,
\]
这是状态函数之间的一个关系，对任意无穷小
变化均成立。由此，对 $n$ 摩尔理想气体，
\[
S(T, V) = nC_{V,m}\ln T + nR\ln V + \text{常数}, \qquad
S(T, P) = nC_{P,m}\ln T - nR\ln P + \text{常数},
\]
对热容为 $C$ 的凝聚相，$S(T) = C\ln T + \text{常数}$。
\end{proposition}

\begin{proof}
可逆时，$T_{\mathrm{ext}} = T$ 且 $S_{\mathrm{created}} = 0$，于是 $\dd S =
\delta Q/T$。随后由第一定律得到恒等式，它不再
依赖于路径。理想气体：$\dd S = (\dd U + P\,\dd V)/T = nC_{V,m}\dd T/T
+ nR\,\dd V/V$；积分；对另一形式利用 $PV = nRT$。凝聚
相：$\dd V = 0$，$\dd S = C\,\dd T/T$。
\end{proof}

\begin{example}[读懂这些公式]\label{ex:b1:second-law-entropy:formulas}
绝热可逆过程是\emph{等熵的}：由 $S(T, V)$
为常数得 $T^{C_{V,m}}V^R$ 为常数，即\ $TV^{\gamma - 1}$ 为常数——
拉普拉斯定律再次出现。一摩尔气体在
恒温下体积加倍，熵增加 $R\ln2 = \qty{5.8}{J/K}$，无论是
可逆进行（熵由热源以热量形式流入），或
通过向真空的焦耳自由膨胀进行（无热量：此时同样的 \qty{5.8}{J/K} 则
被\emph{产生}）。状态函数并不关心它如何到达此状态；
而熵平衡则会计入这一过程。
\end{example}

\section{熵平衡}

\begin{proposition}[与热源交换的热量]\label{prop:b1:second-law-entropy:thermostat}
系统接收来自热源的热量 $Q$，热源温度为 $T_0$（无论
系统内部发生何种变化）有 $S_{\mathrm{exch}} = Q/T_0$，而热源自身的
熵变为 $-Q/T_0$：热源以可逆方式交换熵。
\end{proposition}

\begin{proof}
$T_{\mathrm{ext}} = T_0$ 为常数；热源的内部状态
发生无穷小变化，因此其自身产生的熵可以忽略。
\end{proof}

\begin{example}[冷湖中的热物体]\label{ex:b1:second-law-entropy:quench}
一块铁块（$m = \qty{1.0}{kg}$，$c = \qty{450}{J/(kg.K)}$）温度为 $T_1 =
\qty{373}{K}$，投入温度为 $T_0 = \qty{293}{K}$ 的湖中，最终达到 $T_0$，
铁块的熵变为 $\Delta S_{\text{铁块}} = mc\ln(T_0/T_1) = \qty{-109}{J/K}$，而传给湖的热量为 $Q =
mc(T_0 - T_1) = \qty{-36}{kJ}$，于是 $S_{\mathrm{exch}} =
Q/T_0 = \qty{-123}{J/K}$，$S_{\mathrm{created}} = -109 + 123 = \qty{+14}{J/K}$。
铁块的熵减小，湖的熵增加得更多：该过程是
不可逆的，该数值度量其程度。
\end{example}

\begin{proposition}[两个推论]\label{prop:b1:second-law-entropy:clausiuskelvin}
\index{克劳修斯表述}\index{开尔文表述}
\begin{enumerate}
  \item \emph{（克劳修斯）}热量不会自发地从冷
        物体流向较热的物体：热量 $Q > 0$ 由 $T_h$ 传到
        温度为 $T_c < T_h$ 的物体时，产生 $S = Q(1/T_c - 1/T_h) > 0$ 的熵；
        逆向过程则会使熵减少。
  \item \emph{（开尔文）}任何循环机器都不能从单一
        热源吸热并对外做功：一个循环中 $\Delta S = 0 = Q/T_0 + S_{\mathrm{created}}$，
        要求 $Q \leq 0$，从而 $W = -Q \geq 0$——机器只能
        \emph{接受}功并放出热量。
\end{enumerate}
\end{proposition}

\begin{proof}
两者均对以下系统应用熵平衡：“两个热源”或
“循环中的机器”，其中 $S$ 为状态函数，使得循环中 $\Delta S
= 0$。
\end{proof}

\begin{definition}[熵图]\label{def:b1:second-law-entropy:TS}
\index{熵图}
在 $(S, T)$ 平面上，可逆过程为一条曲线，所吸收的热量
即为曲线下方的面积：$Q_{\mathrm{rev}} = \int T\,\dd S$。可逆
等温线为水平线，等熵线（可逆绝热线）
为竖直线；可逆循环所围面积等于所吸收的热量，
因而由热力学第一定律，也等于输出的功
（见\cref{ch:b1:heat-engines}）。
\end{definition}

\begin{omfigure}
\begin{tikzpicture}
\begin{axis}[omaxis, axis x line=bottom, axis y line=left, width=8.4cm, height=5.6cm, xlabel={$S$}, ylabel={$T$},
    xlabel style={at={(axis description cs:0.5,-0.08)}, anchor=north},
    xmin=0, xmax=5, ymin=0, ymax=4.2, xtick={1,3.5}, xticklabels={$S_1$,$S_2$}, ytick={1.2,3.2}, yticklabels={$T_c$,$T_h$}]
  \fill[omExo!20] (axis cs:1,0) rectangle (axis cs:3.5,3.2);
  \draw[omThm, thick] (axis cs:1,3.2) -- (axis cs:3.5,3.2);
  \draw[omDef, thick] (axis cs:3.5,3.2) -- (axis cs:3.5,1.2);
  \draw[omThm, thick] (axis cs:3.5,1.2) -- (axis cs:1,1.2);
  \draw[omDef, thick] (axis cs:1,1.2) -- (axis cs:1,3.2);
  \fill[white] (axis cs:1,0) rectangle (axis cs:3.5,1.2);
  \fill[omExo!35] (axis cs:1,1.2) rectangle (axis cs:3.5,3.2);
  \node[font=\footnotesize] at (axis cs:2.25,2.2) {$W$ 输出};
  \node[omThm, font=\footnotesize, above] at (axis cs:2.25,3.25) {$Q_h = T_h\Delta S$};
  \node[omThm, font=\footnotesize, below] at (axis cs:2.25,1.15) {$\abs{Q_c} = T_c\Delta S$};
\end{axis}
\end{tikzpicture}\qquad
\begin{tikzpicture}
\begin{axis}[omaxis, axis x line=bottom, axis y line=left, width=7.0cm, height=5.6cm, xlabel={$T_1/T_2$}, ylabel={$S_{\mathrm{created}}/mc$},
    xlabel style={at={(axis description cs:0.5,-0.1)}, anchor=north},
    xmin=0, xmax=4.2, ymin=0, ymax=0.6, xtick={1,2,3,4}, ytick={0.2,0.4}]
  \addplot[omDef, thick, domain=0.15:4.2, samples=200] {ln((1+x)^2/(4*x))};
  \node[omDef, font=\footnotesize] at (axis cs:2.6,0.45) {$\ln\dfrac{(T_1 + T_2)^2}{4T_1T_2}$};
\end{axis}
\end{tikzpicture}

{\small 左图：在熵图中，由两条等温线和两条等熵线组成的一个可逆循环
在高温处吸收 $T_h\Delta S$ 的热量，
在低温处放出 $T_c\Delta S$ 的热量，并将差值以功输出——
即下一章的卡诺循环。右图：两个相等的物体温度分别为 $T_1$ 与 $T_2$ 时
相互接触所产生的熵：仅当 $T_1 = T_2$ 时为零，
否则为正。}
\end{omfigure}

\begin{method}[列写熵平衡]\label{met:b1:second-law-entropy:method}
\begin{enumerate}
  \item 由状态函数计算系统的 $\Delta S$
        （见\cref{prop:b1:second-law-entropy:identity}）：仅需始末
        状态即可。
  \item 由实际吸收的热量及热源温度计算 $S_{\mathrm{exch}}$：
        对热源为 $\sum Q_i/T_i$，
        否则为 $\int\delta Q/T_{\mathrm{ext}}$。
  \item 由此得 $S_{\mathrm{created}} = \Delta S - S_{\mathrm{exch}}$，它必须
        $\geq 0$——若不满足，则该过程不可能发生（或计算有误）；
        对可逆过程 $= 0$。
\end{enumerate}
\end{method}

\begin{example}[冷热混合]\label{ex:b1:second-law-entropy:mixing}
质量均为 $m$ 的水，温度分别为 $T_1$ 与 $T_2$，在保温瓶中混合，
末温为 $(T_1 + T_2)/2$（由热力学第一定律）；$\Delta S = mc[\ln\frac{T_f}{T_1} + \ln\frac{T_f}{T_2}]
= mc\ln\frac{(T_1 + T_2)^2}{4T_1T_2} \geq 0$（算术平均大于
几何平均），无热量交换，故全部为产生的熵。对 \qty{1}{kg} 的水，
分别取 \qty{300}{K} 与 \qty{400}{K} 时为 \qty{86}{J/K}。周末问题表明，若以
\emph{可逆}方式混合二者——借助理想热机——则
可回收功，末温为 $\sqrt{T_1T_2}$。
\end{example}

\section{熵计量的是什么}

\begin{proposition}[玻尔兹曼公式]\label{prop:b1:second-law-entropy:boltzmann}
\index{玻尔兹曼公式}\index{微观态}
宏观态（给定 $U$、$V$、$N$）可由数目为
$\Omega$ 的微观构型（\emph{微观态}）实现。其熵为
\[
S = k_B\ln\Omega .
\]
孤立系统向微观态数
远多于其他宏观态的宏观态演化：熵的增长即是
趋向压倒性最可能状态的行进。
\end{proposition}
\admitted

\begin{example}[为什么气体不会自动回去]\label{ex:b1:second-law-entropy:count}
$N$ 个分子可以自由地处于盒子的任一半：共有 $2^N$ 种排列，
且等可能；只有一种排列使全部分子都在左半边。让限制在一半中的气体膨胀，
则 $\Omega$ 变为原来的 $2^N$ 倍，于是 $\Delta S =
Nk_B\ln2 = nR\ln2$ ——恰好是上面的焦耳膨胀结果，只是现在有了明确含义：
分子的可能方式多了 $2^N$ 倍。它们自发回到原状的概率为
$2^{-N}$：十个分子时约为千分之一，
一百个分子时为 $10^{30}$ 分之一，而对一摩尔而言，则是一个带有 $10^{23}$ 个零的数。
热力学第二定律并非禁令，而是一种极度偏斜的概率，
至今从未被观察到失效。
\end{example}

\begin{omfigure}
\begin{tikzpicture}
\begin{axis}[omaxis, axis x line=bottom, axis y line=left, width=9.5cm, height=5.2cm, xlabel={左半边的分子数，$k$}, ylabel={排列数},
    xlabel style={at={(axis description cs:0.5,-0.16)}, anchor=north},
    scaled y ticks=false,
    xmin=-0.7, xmax=20.7, ymin=0, ymax=2.05e5, xtick={0,5,10,15,20}, ytick={1e5,2e5}, yticklabels={$10^5$,$2\times10^5$},
    ybar, bar width=5pt]
  \addplot[fill=omDef!40, draw=omDef] coordinates {(0,1) (1,20) (2,190) (3,1140) (4,4845) (5,15504) (6,38760) (7,77520) (8,125970) (9,167960) (10,184756) (11,167960) (12,125970) (13,77520) (14,38760) (15,15504) (16,4845) (17,1140) (18,190) (19,20) (20,1)};
  \node[font=\footnotesize, omDef] at (axis cs:15.5,1.5e5) {$N = 20$: $\binom{20}{k}$};
  \draw[omThm, -Stealth] (axis cs:0.8,1.2e5) -- (axis cs:0.2,8e3);
  \node[font=\footnotesize, omThm, anchor=west] at (axis cs:0.1,1.35e5) {全部在左边：$1$};
\end{axis}
\end{tikzpicture}

{\small 盒中二十个分子：左半边含 $k$ 个分子的排列数。
均分的情形有 $184756$ 种排列，
全部在左边的状态仅有一种。若有 $10^{23}$ 个分子，峰值将变成尖峰，
其相对宽度为 $10^{-11}$：均匀气体并非定律，
它是人们唯一可能观察到的情形。}
\end{omfigure}

\section{练习}

\begin{exercise}[$\star$]\label{exo:b1:second-law-entropy:1}
求 \qty{1.0}{kg} 的冰在 \qty{0}{\degreeCelsius} 熔化时的熵变
（$L_f = \qty{334}{kJ/kg}$），以及 \qty{1.0}{kg} 的水在
\qty{100}{\degreeCelsius} 沸腾时的熵变（$L_v = \qty{2260}{kJ/kg}$）。为什么第二个熵变
要大得多？
\end{exercise}

\begin{exercise}[$\star$]\label{exo:b1:second-law-entropy:2}
求 \qty{1.0}{kg} 的水从 \qty{20}{} 加热到
\qty{80}{\degreeCelsius} 的熵变。它是否取决于加热的方式？
\end{exercise}

\begin{exercise}[$\star$]\label{exo:b1:second-law-entropy:3}
一摩尔理想气体在温度不变的条件下体积加倍。
求 $\Delta S$；若膨胀是可逆的，求交换的熵与产生的熵；
若为向真空中的焦耳膨胀，求同样的量。
\end{exercise}

\begin{exercise}[$\star$]\label{exo:b1:second-law-entropy:4}
由 $S(T, V)$ 证明，理想气体的可逆绝热过程满足
$TV^{\gamma - 1} = $ 常量。在
\emph{突然}绝热压缩（见\cref{exo:b1:first-law:11}）中，熵变是多少？
利用该练习的数据计算之。
\end{exercise}

\begin{exercise}[$\star\star$]\label{exo:b1:second-law-entropy:5}
一块质量为 \qty{1.0}{kg} 的铁块（$c = \qty{450}{J/(kg.K)}$）温度为 \qty{100}{\degreeCelsius}，
被投入温度为 \qty{20}{\degreeCelsius} 的湖中。求铁块的熵变，
湖水的熵变以及宇宙的熵变。
\end{exercise}

\begin{exercise}[$\star\star$]\label{exo:b1:second-law-entropy:6}
两份质量均为 \qty{1.0}{kg} 的水，温度分别为 \qty{300}{K} 和 \qty{400}{K}，
经绝热混合后求末温与产生的熵。试一般地证明
产生的熵 $mc\ln[(T_1 + T_2)^2/4T_1T_2]$ 非负。
\end{exercise}

\begin{exercise}[$\star\star$]\label{exo:b1:second-law-entropy:7}
一份 \qty{1.0}{kJ} 的热量由温度为 \qty{300}{K} 的房间经墙壁漏向
温度为 \qty{280}{K} 的室外，求产生的熵。若同样的热量来自
温度为 \qty{1000}{K} 的炉子而传入房间，请比较两者，并讨论“损失”
在何处最大。
\end{exercise}

\begin{exercise}[$\star\star$]\label{exo:b1:second-law-entropy:8}
由一台与单一热源接触并运行循环的机器的熵平衡，证明开尔文表述。
为什么轮船的热机不能
简单地从海水中吸取热量？
\end{exercise}

\begin{exercise}[$\star\star$]\label{exo:b1:second-law-entropy:9}
一摩尔氮气与一摩尔氧气，温度 $T$ 与压强 $P$ 相同，
被隔板分隔在总体积为 $2V$ 的盒中，允许两者混合。求熵变
（将每种气体视为膨胀至整个盒子）。它是产生的还是交换的？
若隔板两侧都是氮气，结果如何？
\end{exercise}

\begin{exercise}[$\star\star\star$]\label{exo:b1:second-law-entropy:10}
一只 \qty{100}{W} 的灯泡工作一小时，其全部能量最终以热量形式进入
温度为 \qty{293}{K} 的房间。求产生的熵，
并从物理上说明它是在何处产生的？
\end{exercise}

\begin{exercise}[$\star\star\star$]\label{exo:b1:second-law-entropy:11}
一个热容为 $C$ 的物体温度为 $T_1$，现将其冷却至 $T_0 < T_1$：
（a）与温度为 $T_0$ 的热源直接接触；（b）经由温度为 $\sqrt{T_1T_0}$ 的中间热源分两步进行；
（c）经由连续的一系列热源。
求每种情形下产生的熵，
并由此总结“可逆冷却”需要满足什么条件。
\end{exercise}

\begin{exercise}[$\star\star\star$]\label{exo:b1:second-law-entropy:12}
求 $N$ 个分子全部位于盒子左半边的概率，其中
$N = 10$、$100$、$10^{22}$（最后一个请写成 10 的幂）。证明
$\Delta S = k_B\ln(2^N)$ 等于 $nR\ln2$，
并评论“不可逆”与“不可能”这两个词。
\end{exercise}

\section{问题：咖啡杯与时间之箭}

\begin{problem}[{周末问题——一杯冷却的咖啡杯、一个混合浴、一台冰箱与一盒分子：
四个熵的平衡，以及所产生的熵若折合为功，
其值为多少}]
\label{pb:b1:second-law-entropy:1}
水：$c = \qty{4180}{J/(kg.K)}$；房间（热源）：$T_0 = \qty{293}{K}$；
$R = \qty{8.314}{J/(mol.K)}$, $k_B = \qty{1.38e-23}{J/K}$.

\medskip
\noindent\textbf{第一部分——冷却的杯子。}
一杯 \qty{0.25}{kg} 的咖啡，温度为 $T_1 = \qty{363}{K}$，冷却至
室温。
\begin{enumerate}
  \item 传给房间的热量，以及房间的熵变。
  \item 咖啡的熵变。
  \item 所产生的熵，并检验其符号。
  \item 杯中咖啡降到 \qty{330}{K} 时被饮用：
        截至该时刻所产生的熵。
  \item 证明对任意 $T_1 > T_0$，$S_{\mathrm{created}} = mc[T_1/T_0 - 1 -
        \ln(T_1/T_0)] \geq 0$（可考察函数 $f(x) = x - 1 - \ln x$）。
  \item 搅拌或对其吹气时咖啡冷却更快：这是否改变
        所产生的熵？理由何在？
\end{enumerate}

\medskip
\noindent\textbf{第二部分——混合：不当与得当。}
两份质量各为 \qty{1.0}{kg} 的水，温度分别为 $T_1 = \qty{400}{K}$ 与 $T_2 =
\qty{300}{K}$，装于一个保温瓶中。
\begin{enumerate}[resume]
  \item 直接混合：求终温 $T_f$ 及所产生的熵。
  \item 改用理想热机，从热水吸热、向冷水放热并输出功，
        直至二者温度相同，
        温度均为 $T_f'$；整个过程可逆。写出熵的平衡式，
        并证明 $T_f' = \sqrt{T_1T_2}$。
  \item 计算 $T_f'$ 与所输出的功（对两份水运用
        热力学第一定律）。
  \item 数值检验可逆方案的熵平衡：
        热水失去的熵等于冷水获得的熵。
  \item 将所作的功与第 6 题中的 $T_0 \times S_{\mathrm{created}}$ 相比较
        （此处以房间温度为参考温度，$T_0 \approx T_2$）：
        粗心混合所产生的熵测量的是什么？
  \item 为什么没有任何厨房能回收这份功？
\end{enumerate}

\medskip
\noindent\textbf{第三部分——冰箱。}
一台冰箱使其内部保持在 $T_c = \qty{278}{K}$，置于
温度为 $T_0$ 的房间中；热量以 \qty{40}{W} 的速率漏入箱内，由机器抽回。
\begin{enumerate}[resume]
  \item 在一秒内写出该机器（循环装置）的熵平衡式：
        从冷的内部接收的熵、传给房间的熵，
        以及内部产生的熵。
  \item 由此推出每秒所需的最小功（即电功率），
        以及相应的最大性能系数
        $Q_c/W$.
  \item 仅用 $T_c$ 和 $T_0$ 表示该上界，并对冷冻柜
        （$T_c = \qty{255}{K}$）求其数值。
  \item 同等大小的实际冰箱耗电约 \qty{15}{W}，
        求其每秒产生的熵。
  \item 若冰箱门敞开，热量将以 \qty{200}{W} 的功率进入：
        求此时的最小功率，并说明敞开的冰箱为何会使房间变暖。
  \item 对于一台完全不需要电能的冰箱，
        热力学第二定律会作何评论？
\end{enumerate}

\medskip
\noindent\textbf{第四部分——计数。}
一只盒子被一堵隔墙等分为两半；理想气体的 $N$ 个分子
位于左半边。撤去隔墙。
\begin{enumerate}[resume]
  \item 分别求前后的微观态数（每个分子可位于
        任一半边），并由玻尔兹曼公式求 $\Delta S$。
  \item 与焦耳自由膨胀的热力学公式所给出的
        $\Delta S = nR\ln2$ 作比较。
  \item 对 $N = 50$，求所有分子在某一瞬时重新全部位于左半边的概率；
        对于一摩尔气体，该概率的数量级
        以 10 的幂次表示。
  \item 需要多少个是非问题（比特）来确定
        $N$ 个分子各自位于哪一半？并与 $\Delta S/(k_B\ln2)$ 比较。
  \item 焦耳自由膨胀中交换与产生的熵各为多少；为什么它是
        不可逆性的教科书式典型例子？
  \item 将同一气体可逆地、等温地压缩回左半边，
        求所接受的功、传给房间的热量与交换的熵；
        为抵消自由膨胀，需外界提供什么？
  \item 用三句话概括：热力学第二定律守恒什么
        （即什么也不守恒）、它禁止什么，以及所产生的熵付出了怎样的代价。
\end{enumerate}
\end{problem}
